\cotefile{144}
\scanpage{62}
\mspage{6}

C'est ainsi que les \inserted{\unsure{cat. des}} cartes cellulaires triangulées \unsure{pondérées} \unsure{apparaissent}\deleted{\unsure{comme}} comme équivalentes à celles des \inserted{surfaces} revêtements \inserted{ramifiés} de \unsure{$\ast$} $\mathbb P^1(\mathbb C)$, ramifiés au plus en les points $0,1,\infty$. Ces revêtements sont de façon canonique des courbes algébriques (propres et lisses) sur $\mathbb C$, et la \deleted{carte} \inserted{donnée $K$ sur $X$} est l'image inverse de \inserted{l'équateur} $\mathbb P^1(\mathbb R)$, les sommets du type $i$ \deleted{sont} \inserted{forment} la fibre de $X$ au dessus du point $i\in\{0,1,\infty\}\subset\mathbb P^1(\mathbb C)$, les ordres réduits des sommets sont les \unsure{pts} de ramification; les arêtes sont les \unsure{comp.} connexes des images inverses des segments ouverts $]1,\infty[$, $]\infty,0[$, $]0,1[$.

Le cas des \inserted{(subdivision barycentrique des)} cartes ordinaires est celui où on a
\[
\rho_1^2=1,
\]
cette dernière se récupère en termes du revêtement en prenant l'image inverse du segment $[0,1]\subset U_{0,3}\subset\mathbb C$, les sommets (ordinaires) sont les pts au dessus de $0$, tandis que les points au dessus de $1$ (resp. $\infty$) sont par définition les centres des arêtes resp. centres des faces. On trouve la carte \inserted{orientée} ``universelle'' \deleted{orientée} sur $\mathbb P^1(\mathbb C)$, comme \deleted{le segment},
\[
S_0=\{0\}\subset K_\infty=[0,1]
\]
\margin{\unsure{$\ast$} (21)}
\margin{61}
